% -------------------------------------------------------------------------
% WBS ID: RES-PHY-LOD
% Deliverable: Hydrodynamic Loading and Impact Physics
% Status: In progress
% Evidence status: Local hydrodynamic loading observables established
% Review status: Not reviewed
% Last updated: 2026-07-19
% Dependencies: RES-MOD-EQS, RES-MOD-SRC, RES-VER-MET
% Downstream interfaces: RES-NUM-SPEC, RES-STR-LOD, RES-VER-MET
% -------------------------------------------------------------------------

\section{RES-PHY-LOD --- Hydrodynamic Loading and Impact Physics}
\label{sec:res-phy-lod}

\subsection{Purpose and Scope}
\label{sec:res-phy-lod-purpose}

% TODO: Define what this Deliverable covers and explicitly excludes.

This Deliverable defines the hydrodynamic loading observables that the
local three-dimensional URANS--VOF model must resolve for barrier
assessment and later structural handoff. The active baseline assumes a
fixed rigid barrier; structural deformation, material damage, scour and
two-way FSI are stretch extensions after hydrodynamic load extraction
is stable.

\subsection{Research Questions}
\label{sec:res-phy-lod-research-questions}

% TODO: State the questions that must be answered before the Deliverable
% can be accepted.

\subsection{Terminology and Definitions}
\label{sec:res-phy-lod-definitions}

% TODO: Define the technical terminology, symbols, quantities and
% classifications used in this Deliverable.

\subsection{Literature Evidence}
\label{sec:res-phy-lod-literature-evidence}

% TODO: Record evidence from peer-reviewed papers, authoritative datasets,
% standards, technical reports and primary documentation.
%
% Do not create a detached literature-summary list. Explain how each source
% supports, contradicts or limits a project decision.

\textcite{qin2018comparison} is retained as evidence that local
three-dimensional modelling enables direct structural-force extraction
from pressure and shear traction. \textcite{watanabe2022validation} is
retained as evidence that local force, pressure and velocity validation
must be assessed separately from water-level agreement. Source roles:
`3D-RANSVOF` and `3D-VALIDATION`. Project conclusion: the local model
shall output complete pressure and traction histories, not only
empirical drag coefficients. Limitation: these quantities remain
hydrodynamic loads until structural response or FSI is activated.

\subsection{Governing Relations and Quantitative Definitions}
\label{sec:res-phy-lod-governing-relations}

% TODO: Record relevant equations, dimensionless groups, empirical models,
% extraction rules or quantitative definitions.
%
% Use align environments for displayed equations.
% Every displayed equation must have a unique \label{eq:...}.
% Do not place punctuation after displayed equations.

The pressure and viscous contributions to traction on a barrier
surface are:

\begin{align}
  \mathbf{t}_{p}
  &=
  -p\mathbf{n}
  \label{eq:res-phy-lod-pressure-traction}
  \\
  \mathbf{t}_{\tau}
  &=
  \boldsymbol{\tau}
  \mathbf{n}
  \label{eq:res-phy-lod-shear-traction}
  \\
  \mathbf{t}
  &=
  \mathbf{t}_{p}
  +
  \mathbf{t}_{\tau}
  \label{eq:res-phy-lod-total-traction}
\end{align}

The loading quantities of interest are:

\begin{align}
  \mathbf{F}(t)
  &=
  \int_{\Gamma_{\mathrm{B}}}
  \mathbf{t}
  \,\dd\Gamma
  \label{eq:res-phy-lod-force}
  \\
  \mathbf{M}(t)
  &=
  \int_{\Gamma_{\mathrm{B}}}
  \left(
    \mathbf{x}-\mathbf{x}_0
  \right)
  \times
  \mathbf{t}
  \,\dd\Gamma
  \label{eq:res-phy-lod-moment}
  \\
  \mathbf{J}
  &=
  \int_{t_0}^{t_1}
  \mathbf{F}(t)
  \,\dd t
  \label{eq:res-phy-lod-force-impulse}
\end{align}

\subsection{Physical Mechanisms}
\label{sec:res-phy-lod-physical-mechanisms}

% TODO: Complete this RES-PHY Domain-specific subsection.

The local impact model shall represent hydrostatic loading,
nonhydrostatic impact pressure, bore impact, reflected and transmitted
wave loading, overtopping-induced loading, shear traction, wake effects
and pressure-duration effects. Structural response is not inferred from
these hydrodynamic quantities unless `RES-STR` activates a structural
model.

\subsection{Characteristic Spatial and Temporal Scales}
\label{sec:res-phy-lod-characteristic-spatial-temporal-scales}

% TODO: Complete this RES-PHY Domain-specific subsection.

\subsection{Relevant Observables}
\label{sec:res-phy-lod-relevant-observables}

% TODO: Complete this RES-PHY Domain-specific subsection.

Required local loading observables are pressure histories, pressure
spatial distributions, shear traction, resultant force, force impulse,
overturning moment, torsional moment, centre-of-pressure migration,
impact duration and repeat-wave cumulative loading.

Each loading observable shall carry the identifiers
\texttt{corridor\_id}, \texttt{geometry\_id}, \texttt{patch\_id},
local coordinate frame, integration interval and accepted-timestep
status. Source role: \textcite{qin2018comparison} supports direct
force extraction from local 3D fields; \textcite{watanabe2022validation}
supports reporting pressure, velocity and force separately. Downstream
handoff: RES-STR-LOD shall consume these histories as hydrodynamic
loads, not as proof of material response.

\subsection{Controlling Dimensionless Parameters}
\label{sec:res-phy-lod-controlling-dimensionless-parameters}

% TODO: Complete this RES-PHY Domain-specific subsection.

\subsection{Physical Regimes and Transition Conditions}
\label{sec:res-phy-lod-physical-regimes-transition-conditions}

% TODO: Complete this RES-PHY Domain-specific subsection.

\subsection{Implications for Model Validity}
\label{sec:res-phy-lod-model-validity-implications}

% TODO: Complete this RES-PHY Domain-specific subsection.

\subsection{Candidate Approaches}
\label{sec:res-phy-lod-candidate-approaches}

% TODO: Define the credible candidate approaches considered.

\subsection{Critical Comparison}
\label{sec:res-phy-lod-critical-comparison}

% TODO: Compare candidates using physically and project-relevant criteria.
% Do not select an approach solely on popularity or implementation ease.

\subsection{Assumptions and Applicability}
\label{sec:res-phy-lod-assumptions}

% TODO: State assumptions and the conditions under which they are valid.

\subsection{Limitations and Failure Conditions}
\label{sec:res-phy-lod-limitations}

% TODO: State where the evidence, model or selected approach ceases to be
% reliable or applicable.

\subsection{Project-Specific Conclusions}
\label{sec:res-phy-lod-conclusions}

% TODO: Convert the evidence into conclusions specific to the Studentship
% project. Do not merely restate literature findings.

\subsection{Selected Approach and Rationale}
\label{sec:res-phy-lod-selected-approach}

% TODO: Record the selected approach only after sufficient evidence exists.
% State the alternatives rejected and the reasons for rejection.

The selected loading approach is direct traction integration over the
resolved barrier surface. Empirical drag-only loading may be retained
for screening comparisons or sanity checks but shall not replace the
local pressure and shear extraction used for validation and structural
handoff.

\subsection{Unresolved Decisions}
\label{sec:res-phy-lod-unresolved-decisions}

% TODO: Record unresolved questions, required evidence and decision owners.

\subsection{Requirements and Handoffs}
\label{sec:res-phy-lod-requirements}

% TODO: State requirements passed to other Research Domains, Software,
% verification activities or later execution work.

`RES-NUM-SPEC` shall extract the loading metrics after each accepted
local timestep. `RES-VER-MET` shall define acceptance criteria for
pressure, force, impulse and moment. `RES-STR-LOD` shall decide how
accepted hydrodynamic tractions are transferred into structural load
cases. Software shall emit pressure, shear, force, moment and impulse
histories in the structured output schema before any structural damage
or FSI extension is attempted.

\subsection{Deliverable Completion Record}
\label{sec:res-phy-lod-completion-record}

% TODO: Record whether the Deliverable is:
%
% - Not started
% - In progress
% - Draft complete
% - Reviewed
% - Accepted
%
% Acceptance requires:
%
% - the research questions to be answered;
% - claims to be supported by appropriate evidence;
% - assumptions and limitations to be explicit;
% - the project-specific conclusion to be stated;
% - downstream requirements to be traceable;
% - unresolved decisions to be closed or formally deferred.
